% PFC: PhotoPlace
% (c) Jose Riguera Lopez 2013

% Opciones:
%       a4paper -> indica el tamaño del papel, en este caso A4.
%       11pt    -> tamaño de la fuente 11 puntos.
%       twoside -> doble cara.
%       oneside -> una cara del folio.
%
\documentclass[a4paper,11pt,twoside]{book}
%% Estilo PFC UDC
\usepackage{pfc}
%% Macros
\newcommand{\curso}{Ingeniería en Informática}
\newcommand{\proyecto}{Proyecto Fin de Carrera}
\newcommand{\universidad}{Universidad de A Coruña}
\newcommand{\centro}{Facultad de Informática}
\newcommand{\unicentro}{\centro de la \universidad}
\newcommand{\proyectotipo}{Proyecto clásico de ingeniería}
\newcommand{\titulo}{Cómo el pollo cruzó la carretera}
\newcommand{\titulogalego}{Título longo do proxecto}
\newcommand{\tituloenglish}{Long title of the project}
\newcommand{\corto}{Título corto}
\newcommand{\autor}{Dolores Fuertes de Barriga}
\newcommand{\director}{Armando Bronca Segura}
\newcommand{\tutor}{Alex Plosivo Seguro}
\newcommand{\departamento}{Departamento de asuntos humorísticos}
\newcommand{\software}{PhotoPlace}
\newcommand{\palabras}{lista, palabras, clave}

\title{\titulo}
\author{\autor}
\date{\today}

\hypersetup{
%   bookmarks=true,         % barra de marcadores
    unicode=false,          % caracteres non-Latin en marcadores de Acrobat
    pdftoolbar=true,        % mostrar barra de herramientas de Acrobat
    pdfmenubar=true,        % mostrar menú de Acrobat
    pdffitwindow=false,     % ajustar ventana al ancho de página
    pdfstartview={FitH},    % ajustar documento al ancho de página
    pdftitle={\corto},      % título
    pdfauthor={\autor},     % autor
    pdfsubject={\titulo},   % tema del documento
    pdfcreator={LaTex},     % generador del documento
    pdfproducer={\autor},   % productor del documento
    pdfkeywords={\palabras},% lista de palabras clave
    pdfnewwindow=true,      % enlaces en una nueva ventana
    colorlinks=true,        % false: enlaces en caja; true: enlaces coloreados
    linkcolor=black,        % color de enlaces internos
    citecolor=green,        % color de enlaces a bibliografía
    filecolor=magenta,      % color de enlaces a ficheros
    urlcolor=blue           % color de enlaces externos
}

\loadglsentries{pfc.glossary}
\makeglossaries

\begin{document}
  \graphicspath{{./images/}}
  % pageblank genera una pagina en blanco, para imprimir en formato libro
  %
% Portada.
%
\begin{titlepage}
	\begin{center}
		% Logotipo de la universidad.
		\includegraphics[width=6cm]{./eps/logo_udc.eps}
		\vspace{2cm}

		% Nombre de la facultad, de la universidad y del departamento.
		{\Large{\textbf{\centro}}}
		\\
		{\it \large{\textbf{\departamento}}}
		\vspace{1cm}

		{\large {\sc \proyecto}\\{\curso}}
		\vspace{1cm}

		% Título
		\textbf{\Large \titulo}
		\vspace{6cm}
	\end{center}

	\begin{flushright}
		\begin{tabular}{ll}
			\large{\textbf{Alumno:}}	&
			\large{\autor} \\

			\large{\textbf{Director:}}	&
			\large{\director} \\

			\large{\textbf{Tutor:}}		&
			\large{\tutor} \\

			% Fecha.
			\large{\textbf{Fecha:}}		&
			\large{\today} \\
		\end{tabular}
	\end{flushright}
\end{titlepage}

  \pageblank
  %
% Certificado
%
\thispagestyle{plain}
\begin{center}
	\begin{minipage}[t][6cm][l]{.8\textwidth}
		\begin{center}
			% Director del proyecto
			D. {\sc \director}

			% Profesores
			Profesor, \centro

			% Departamento al que pertenece el director y en el que se realiza el proyecto.
			\departamento
			
			\universidad
		\end{center}
	\end{minipage}
\end{center}

CERTIFICA:

Que la memoria titulada {\it ``\titulo''} ha sido realizada por {\sc \autor}
de acuerdo a la descripción inicialmente propuesta, bajo mi dirección y
constituye su {\proyecto} de {\curso}.
Por la presente, autorizo su presentación para que el Proyecto sea defendido
en esta convocatoria.

\vspace{3cm}

En A Coruña, a \today

% Espacio para que pueda firmar el certificado que debe acompañar al proyecto.
\vspace{3cm}
\begin{center}
	\begin{minipage}[t][4cm][l]{.5\textwidth}
	D. {\sc \director}
	\\
	Director del proyecto
	\end{minipage}
\end{center}

  \pageblank
  %
% Nota
%
\thispagestyle{plain}
\section*{}


\begin{tabular}{p{2cm}p{11cm}}
	\large{Título} & \\
	& \textbf{\large{\titulo}} \\
	\\
	& \textbf{\large{\titulogalego}} \\
	\\
	& \textbf{\large{\tituloenglish}} \\
	\\
	\large{\textbf{Clase:}} & \large{\proyectotipo} \\
	\\
	\large{\textbf{Autor:}} & \large{\autor} \\
	\large{\textbf{Tutor:}} & \large{\tutor} \\
	\large{\textbf{Director:}} & \large{\director} \\
	\large{\textbf{Fecha:}} & \large{\today} \\
	\\
	\large{\textbf{Tribunal}} & \\
	& \vspace{3cm} \\
	\\
	\large{\textbf{Fecha de defensa:}} & \\
	\\
	\large{\textbf{Calificación:}} & \\
	\\
\end{tabular}

  \pageblank
  \thispagestyle{plain}
\section*{Resumen}

A lo largo de los últimos años ...

  \pageblank
  \thispagestyle{plain}
\section*{Lista de Palabras Clave}

{\setlength{\parindent}{0cm}\ttfamily
\palabras
}



  \pageblank
  %
% Dedicatoria
%
\thispagestyle{plain}

\begin{flushright}
	{\it A todo el mundo.}
\end{flushright}

  \pageblank
  %
% Agradecimientos
%
\chapter*{Agradecimientos}
\vspace*{2cm}

A los profesores {\director} y {\tutor} por sus consejos durante el desarrollo del proyecto. \\

A María de la O por ... \\

  \pageblank
  
  % FRONTMATTER: TOC, LOF, LOT
  \frontmatter
  \dominitoc[n]     % removes the title "contents"
  \nomtcrule        % removes rules = horizontal lines
  \tableofcontents
  %\addcontentsline{toc}{chapter}{Contents}
  \listoffigures
  %\addcontentsline{toc}{chapter}{List of Figures}
  %\listoftables
  %\addcontentsline{toc}{chapter}{List of Tables}

  % MAINMATTER: content
  \mainmatter
  \chapter[Introducción]{
  \label{chp:introduccion}
  Introducción
}
\minitoc
\newpage

El proceso de desarrollo del software comienza con un análisis del entorno y
del dominio de la aplicación. En un contexto amplio, se entiende por dominio
como ``una esfera de actividad o interés'', también denominado campo. Es decir,
el área de aplicación donde el software fue creado.

En el presente capítulo se describe el entorno y el dominio de la aplicación, y la
evolución de los servicios de geolocalización en los últimos años. A su vez,
también se exponen los objetivos de este proyecto, dando una visión
de su estructura y las partes que la forman.


\section{Dominio de Aplicación}

La evolución de los \glspl{SIG} ha sido imparable en la ultima década. Motivado 
-sin duda- por el interés de gigantes tecnológicos como Google, Apple y Microsoft, 
ha provocado que se hayan popularizado y extendido tecnologías que  antes sólo 
estaban al alcance de profesionales y científicos.


  \chapter[Conceptos Previos]{
  \label{chp:conceptos}
  Conceptos Previos
}
\minitoc
\newpage

En los últimos años, la información geográfica -también llamada información
geoespacial- a cobrado una gran importancia. La facilidad -y necesidad- de
identificar el lugar donde se han obtenido los datos ha provocado la
aparición de nuevos términos. A continuación se presenta una serie de
definiciones que ayudarán a comprender el dominio y alcance de este proyecto.


\section{Coordenadas Geográficas}

Para determinar la posición de un objeto en el globo terrestre se emplean
coordenadas geográficas, un sistema de referencia que utiliza  dos
coordenadas angulares: latitud (eje Norte-Sur) y longitud (eje Este-Oeste)
y una tercera coordenada elevación para indicar la altitud.

\begin{description}
 \item[Latitud] es el ángulo que existe entre un punto cualquiera y el
 Ecuador, medida sobre el meridiano que pasa por dicho punto. En cartografía,
 la latitud se suele expresar en grados sexagesimales. Todos los puntos
 ubicados sobre el mismo paralelo tienen la misma latitud. Aquellos que se
 encuentran al norte del Ecuador reciben la denominación Norte (N), los que se
 encuentran al sur del Ecuador reciben la denominación Sur (S).
 El rango de valores es de 0º para el Ecuador a 90º para los polos Norte y Sur
 con latitud 90º N y 90º S respectivamente.
 \item[Longitud] es el ángulo a lo largo del ecuador desde cualquier punto de
 la Tierra. Por convenio, está aceptado internacionalmente que Greenwich en Londres
 es la longitud 0. Las líneas de longitud son círculos máximos que pasan por los
 polos y se llaman meridianos. La longitud también se expresa en grados
 sexagesimales entre 0º y 180º indicando a qué hemisferio (occidental W —del
 inglés \textit{West}— y oriental E —\textit{East}—) o también entre 0º y 180º positivos
 indicando Este, o negativos indicando hacia el Oeste.
 \item[Elevación] o \textbf{altitud} es la distancia vertical de un punto de la Tierra
 respecto al nivel del mar, llamada elevación sobre el nivel medio del mar, en contraste
 con la altura, que indica la distancia vertical existente entre dos puntos de la
 superficie terrestre.
\end{description}

Combinando estos dos ángulos (latitud y longitud) junto con la elevación, es posible
expresar la posición de cualquier punto de la superficie de la Tierra.

  \chapter[Plan de Proyecto]{
  \label{chp:plan}
  Plan de Proyecto
}
\minitoc
\newpage

Como paso previo a la realización de un proyecto es necesario realizar una serie de
análisis que -dependiendo del tipo de sistema a desarrollar- determinen su
viabilidad y permitan intentar predecir los problemas y los riesgos que se
encontrarán durante su elaboración. Por ejemplo, en un proyecto de ingeniería,
si los problemas de desarrollo superan a las expectativas vinculadas a las
predicciones y previsiones, no compensaría el esfuerzo de implementación.
En este capítulo, se estudia el plan de proyecto para la propuesta realizada,
sin tener en cuenta las futuras ampliaciones.


\section{Análisis de Viabilidad}

Al emprender el desarrollo de un proyecto los recursos y el tiempo deben ser
realistas para su materialización a fin de evitar pérdidas económicas o
frustración profesional. La viabilidad y el análisis de riesgos están relacionados
de muchas maneras, si el riesgo del proyecto es alto, la viabilidad de producir
software de calidad se reduce. Así pues, en el presente capítulo aborda los siguientes
puntos:

\begin{itemize}
    \item \textbf{Descripción} formal del proyecto.
    \item \textbf{Objeto}, motivos y necesidades del usuario.
    \item \textbf{Objetivos}, declaración del objetivo final del proyecto.
    \item \textbf{Metodología}, introducción al proceso del ciclo de vida elegido.
    \item \textbf{Requerimientos y recursos necesarios} para alcanzar el objetivo.
    \item \textbf{Sistema de control} para el seguimiento de los hitos principales.
    \item \textbf{Elementos de riesgo} que pueden hacer inviable el proyecto.
    \item \textbf{Beneficios} esperados tras la realización del proyecto.
    \item \textbf{Conclusiones} y resoluciones finales.
\end{itemize}

  \chapter[Tecnología]{
  \label{chp:tecnologia}
  Tecnología
}
\minitoc
\newpage

La base tecnológica es el grupo de tecnologías sobre las que se apoya el
software y su desarrollo. Constituyen una parte importante en el análisis
del proyecto, puesto que, de su elección dependen otros factores tales como:
curva de aprendizaje y facilidad de uso, soporte y evolución futura,
documentación y comunidad, etc.

Este capítulo se centra en la última versión, a fecha de entrega de esta memoria,
del programa. Es probable que algunas módulos aquí expuestos sean sustituidos
por otros en futuras \glspl{iteracion} siguiendo el ciclo de vida del software.


\section{Base Tecnológica}

A continuación se describen las tecnologías más importantes sobre las que
se ha desarrollado el programa \textbf{\software}. El nexo común entre todas las 
tecnologías expuestas es su capacidad multiplataforma y formar parte del ecosistema del
software libre.


  \chapter[Diseño]{
  \label{chp:disenho}
  Diseño
}
\minitoc
\newpage

Como se ha visto, la realización del proyecto ha seguido directrices marcadas
por la \gls{metodologia agil} \gls{Scrum-Ban} y por el \gls{ddd}. Además,
la arquitectura de la aplicación se ha dividido en distintas capas apoyadas por
patrones de diseño e integración.

En este capítulo se muestran, en primer lugar, los conceptos básicos necesarios
para comprender el proceso de diseño del software y, posteriormente, los modelos
del software logrados tras la última iteración de la aplicación.


\section{Ciclo de vida}

El ciclo de vida está basado en los principios de las \glspl{metodologia agil} y
se organiza en \glspl{iteracion} que incluyen: análisis de requisitos,
diseño, codificación, revisión y documentación.

  \chapter[Implementación]{
  \label{chp:implementacion}
  Implementación
}
\minitoc
\newpage

En el presente capítulo se prestará atención a la arquitectura de la aplicación,
es decir, a los bloques de construcción del programa. Los componentes son:
paquetes, librerías, frameworks, \textit{APIs}, etc descritos en términos de
módulos, clases, funciones y algoritmos. En esencia, se trata del software y de
la organización del código fuente. Se mostrarán fragmentos de código del
software implementado con los aspectos más significativos. Esta parte de la
documentación técnica es de gran importancia para comprender el proyecto.


\section{Organización del código fuente}

El código fuente de todo el proyecto, incluida la presente memoria realizada en
\LaTeX, se ha gestionando en Git, usando el flujo de trabajo Git-Flow. Se ha
estructurado el código fuente de la siguiente forma:

  \chapter[Conclusiones]{
  \label{chp:conclusiones}
  Conclusiones
}
\minitoc
\newpage

En este último capítulo del proyecto se presenta una estimación de los costes
de desarrollo totales para, después de las conclusiones, finalizar el capítulo
con unos apuntes acerca de las direcciones de las posibles futuras ampliaciones.


\section{Análisis Económico}

En esta sección se muestra un estudio de los costes de la aplicación basándose
principalmente en el tiempo empleado. Se ha realizado un análisis económico
desde  dos puntos de vista diferentes: costes de desarrollo y retorno de inversión.

  
  \appendix
  \chapter[Apéndice]{
  \label{chp:apendice}
  Demostración
}

En este apéndice se mostrará el funcionamiento de la aplicación, primero con
la versión para computador y a continuación la aplicación web. Para ello, se
mostrará el caso de uso para generar una presentación KML. No se
pretende detallar todos los parámetros del proceso, puesto que están cubiertos 
por la documentación de usuario disponible en la web, sólo mostrar la 
apariencia de la aplicación y su filosofía de funcionamiento.

\section{\software}

Como se ha comentado, el proyecto ha sido desarrollado en abierto usando
la plataforma Google Code. El nombre elegido para el software ha sido
\software\footnote{Web oficial del proyecto \url{https://code.google.com/p/photoplace}} 
y se ha desarrollado un logotipo que refuerza la identidad del programa:

\begin{figure}[H]
  \centering
  \includegraphics[scale=0.4,keepaspectratio=true]{photoplace.png}
  \caption*{Web de desarrollo de \software\ en Google Code}
  \label{fig:photoplace}
\end{figure}

  % Glossary
  \printglossary[title=Glosario,toctitle=Glosario]
  \printglossary[type=\acronymtype,title=Acrónimos,toctitle=Acrónimos]
  % Bibliography
  \nocite{*}    % incluir referencias no citadas.
  %\bibliographystyle{plainnat}
  %\bibliographystyle{unsrtnat}
  \bibliographystyle{dcu}  % Books and web
  \bibliography{pfc}
  
  %
% Licencia
%
\pagestyle{plain}

\begin{flushright}
\begin{scriptsize}
Everything will be okay in the end. 
If it's not okay, it's not the end.
\end{scriptsize}
\end{flushright}

\vspace*{17cm}

\begin{flushleft}
\begin{scriptsize}
\begin{verbatim}
Copyright (c) José Riguera.
Permission is granted to copy, distribute and/or modify this
document under the terms of the GNU Free Documentation License,
Version 1.2 or any later version published by the Free Software
Foundation; with the Invariant Sections being: DEDICATORIA,
AGRADECIMIENTOS and BIBLIOGRAFÍA, with the Front-Cover Texts
being PORTADA UDC, and with no Back-Cover Texts.
\end{verbatim}
\end{scriptsize}
\end{flushleft}

\end{document}
